\begin{env}[4.1.5]
\label{0.4.1.5-fr}
Dado un espacio anillado $(X,\mathcal{A})$ (no necesariamente conmutativo),
no recordaremos aquí las definiciones de los bifuntores
$\mathcal{F}\otimes_{\mathcal{A}}\mathcal{G}$,
$\mathcal{H}\!om_{\mathcal{A}}(\mathcal{F},\mathcal{G})$ y
$\operatorname{Hom}_{\mathcal{A}}(\mathcal{F},\mathcal{G})$
(G, II, 2.8 y 2.2) en las categorías de los $\mathcal{A}$-Módulos a la izquierda
o a la derecha (según el caso), con valores en la categoría de los haces de
grupos abelianos (o, en términos más
\oldpage[0\textsubscript{I}]{38}
generales, en la de los $\mathcal{C}$-Módulos, si $\mathcal{C}$ es el centro de
$\mathcal{A}$). El tallo
$(\mathcal{F}\otimes_{\mathcal{A}}\mathcal{G})_x$ en todo punto $x\in X$
se identifica canónicamente con
$\mathcal{F}_x\otimes_{\mathcal{A}_x}\mathcal{G}_x$, y se define un
homomorfismo canónico funtorial
\[
  (\mathcal{H}\!om_{\mathcal{A}}(\mathcal{F},\mathcal{G}))_x
  \longrightarrow
  \operatorname{Hom}_{\mathcal{A}_x}(\mathcal{F}_x,\mathcal{G}_x)
\]
que, en general, no es ni inyectivo ni sobreyectivo. Los bifuntores que
acabamos de considerar son aditivos y, en particular, conmutan con las sumas
directas finitas. El funtor $\mathcal{F}\otimes_{\mathcal{A}}\mathcal{G}$ es
exacto por la derecha en $\mathcal{F}$ y en $\mathcal{G}$, conmuta con los
límites inductivos, y $\mathcal{A}\otimes_{\mathcal{A}}\mathcal{G}$
(resp. $\mathcal{F}\otimes_{\mathcal{A}}\mathcal{A}$) se identifica
canónicamente con $\mathcal{G}$ (resp. $\mathcal{F}$). Los funtores
$\mathcal{H}\!om_{\mathcal{A}}(\mathcal{F},\mathcal{G})$ y
$\operatorname{Hom}_{\mathcal{A}}(\mathcal{F},\mathcal{G})$ son
\emph{exactos por la izquierda} en $\mathcal{F}$ y $\mathcal{G}$.
Más precisamente, si se tiene una sucesión exacta de la forma
$0\to\mathcal{G}'\to\mathcal{G}\to\mathcal{G}''$, la sucesión
\[
  0\longrightarrow
  \mathcal{H}\!om_{\mathcal{A}}(\mathcal{F},\mathcal{G}')
  \longrightarrow
  \mathcal{H}\!om_{\mathcal{A}}(\mathcal{F},\mathcal{G})
  \longrightarrow
  \mathcal{H}\!om_{\mathcal{A}}(\mathcal{F},\mathcal{G}'')
\]
es exacta, y si se tiene una sucesión exacta de la forma
$\mathcal{F}'\to\mathcal{F}\to\mathcal{F}''\to0$, la sucesión
\[
  0\longrightarrow
  \mathcal{H}\!om_{\mathcal{A}}(\mathcal{F}'',\mathcal{G})
  \longrightarrow
  \mathcal{H}\!om_{\mathcal{A}}(\mathcal{F},\mathcal{G})
  \longrightarrow
  \mathcal{H}\!om_{\mathcal{A}}(\mathcal{F}',\mathcal{G})
\]
es exacta. El funtor $\operatorname{Hom}$ tiene propiedades análogas.
Además, $\mathcal{H}\!om_{\mathcal{A}}(\mathcal{A},\mathcal{G})$ se identifica
canónicamente con $\mathcal{G}$. Por último, para todo abierto $U\subset X$,
se tiene
\[
  \Gamma\bigl(U,\mathcal{H}\!om_{\mathcal{A}}(\mathcal{F},\mathcal{G})\bigr)
  =
  \operatorname{Hom}_{\mathcal{A}|U}(\mathcal{F}|U,\mathcal{G}|U).
\]

Para todo $\mathcal{A}$-Módulo a la izquierda $\mathcal{F}$ (resp. a la
derecha), se llama \emph{dual} de $\mathcal{F}$ y se denota por
$\check{\mathcal{F}}$ el $\mathcal{A}$-Módulo a la derecha (resp. a la
izquierda) $\mathcal{H}\!om_{\mathcal{A}}(\mathcal{F},\mathcal{A})$.

Por último, si $\mathcal{A}$ es un haz de anillos conmutativos y
$\mathcal{F}$ es un $\mathcal{A}$-Módulo,
$U\to\bigwedge^p\Gamma(U,\mathcal{F})$ es un prehaz cuyo haz asociado es un
$\mathcal{A}$-Módulo que se denota por
$\bigwedge^p\mathcal{F}$ y se llama la \emph{potencia exterior $p$-ésima
de $\mathcal{F}$}. Se comprueba fácilmente que la aplicación canónica del
prehaz $U\to\bigwedge^p\Gamma(U,\mathcal{F})$ en el haz asociado
$\bigwedge^p\mathcal{F}$ es \emph{inyectiva} y que, para todo
$x\in X$, se tiene
$(\bigwedge^p\mathcal{F})_x=\bigwedge^p(\mathcal{F}_x)$. Es claro que
$\bigwedge^p\mathcal{F}$ es un funtor covariante en $\mathcal{F}$.
\end{env}
